% SALAD: SOC Alert Labeled Analysis Dataset
% Data in Brief --- Dataset Paper

\documentclass[review]{elsarticle}
\usepackage{booktabs}
\usepackage{multirow}
\usepackage{hyperref}
\usepackage{graphicx}
\usepackage{xcolor}
\usepackage{amssymb}
\usepackage{amsmath}

\journal{Data in Brief}

\begin{document}

\begin{frontmatter}

\title{SALAD: SOC Alert Labeled Analysis Dataset --- A Unified Benchmark for Alert Triage, Prioritization, and Classification}

\author[1]{Nutthakorn Chalaemwongwan\corref{cor1}}
\cortext[cor1]{Corresponding author.}
\ead{nutthakorn.ch@kmitl.ac.th}
\affiliation[1]{organization={Department of Computer Engineering, KOSEN-KMITL, King Mongkut's Institute of Technology Ladkrabang}, city={Bangkok}, country={Thailand}}

\begin{abstract}
This paper presents SALAD (SOC Alert Labeled Analysis Dataset), a unified benchmark dataset containing 2,778,424 SOC (Security Operations Center) alerts derived from two established network intrusion detection datasets: CIC-IDS2017 and UNSW-NB15. Each alert record follows a standardized 29-field schema and is enriched with MITRE ATT\&CK tactic and technique mapping, Cyber Kill Chain phase labels, five-level severity scoring, analyst triage decisions (escalate, investigate, or suppress), continuous priority scores, three-level difficulty labels, and natural language alert descriptions. The dataset is split into training (70\%), validation (15\%), and test (15\%) sets with stratified sampling. Four benchmark tasks are defined: binary alert classification, three-class triage prediction, priority score regression, and 15-class attack categorization. Baseline evaluation scripts using Logistic Regression and Random Forest classifiers are provided. The dataset, evaluation code, and data card are publicly available on Zenodo (DOI: 10.5281/zenodo.18912306), HuggingFace, and GitHub under a CC-BY 4.0 license.
\end{abstract}

\begin{keyword}
SOC alert triage \sep MITRE ATT\&CK mapping \sep cybersecurity benchmark \sep network intrusion detection \sep security operations center \sep multi-task evaluation
\end{keyword}

\end{frontmatter}

% ============================================================
\section{Specifications Table}
% ============================================================

\begin{table}[h]
\small
\begin{tabular}{p{4cm} p{9cm}}
\toprule
\textbf{Subject} & Computer Science / Cybersecurity \\
\textbf{Specific subject area} & SOC alert triage, intrusion detection, cybersecurity benchmark \\
\textbf{Type of data} & CSV (tabular), JSON (statistics) \\
\textbf{Data collection} & Derived from CIC-IDS2017 and UNSW-NB15 via an automated five-stage conversion pipeline that adds SOC-specific metadata (triage decisions, priority scores, MITRE ATT\&CK mapping, difficulty levels, and natural language descriptions) not present in the source data \\
\textbf{Data source location} & CIC-IDS2017: Canadian Institute for Cybersecurity, University of New Brunswick, Canada; UNSW-NB15: University of New South Wales, Australia \\
\textbf{Data accessibility} & Repository name: Zenodo, HuggingFace, GitHub \newline DOI: 10.5281/zenodo.18912306 \newline Zenodo: \url{https://doi.org/10.5281/zenodo.18912306} \newline HuggingFace: \url{https://huggingface.co/datasets/nutthakorn7/SALAD-SOC} \newline GitHub: \url{https://github.com/nutthakorn7/SALAD-SOC} \\
\textbf{Related research article} & N/A (standalone data article) \\
\bottomrule
\end{tabular}
\end{table}


% ============================================================
\section{Background}
% ============================================================

Security Operations Centers (SOCs) process thousands to millions of alerts daily from intrusion detection systems, firewalls, and endpoint security tools~\cite{alahmadi2022socsurvey}. Studies report that SOC analysts face an average of 4,484 alerts per day, with 67\% of alerts being ignored due to alert fatigue~\cite{ponemon2023alert}. Efficient alert triage --- deciding which alerts to escalate, investigate, or suppress --- is critical for timely incident response~\cite{husak2019}. Although ML-based~\cite{zuech2015ids_survey} and RL-based approaches~\cite{cage2022cyborg} to alert triage automation are gaining traction, most studies rely on proprietary data or synthetic simulations because no public SOC alert dataset with triage labels exists~\cite{landauer2024,msguide2024}. Existing public datasets such as CIC-IDS2017~\cite{sharafaldin2018} and UNSW-NB15~\cite{moustafa2015} provide raw network traffic but lack SOC-specific metadata: there are no triage decisions, no priority scores, no MITRE ATT\&CK mapping~\cite{strom2018mitre}, and no structured benchmark tasks. SALAD fills this void by transforming these raw traffic records into a realistic SOC alert format, adding five layers of computed metadata that do not exist in the source data.


% ============================================================
\section{Value of the Data}
% ============================================================

\begin{itemize}
\item SALAD provides real network traffic transformed into a unified SOC alert format combining MITRE ATT\&CK mapping, triage decision labels, priority scores, difficulty levels, and natural language descriptions --- properties not simultaneously available in existing public datasets (see Table~\ref{tab:comparison}).

\item The dataset enables researchers across multiple communities --- machine learning, reinforcement learning (RL/MARL), and large language models (LLMs) --- to develop and benchmark alert triage systems using a single, standardized evaluation framework~\cite{veeramachaneni2016aianalyst}.

\item With 2,778,424 alerts across 15 attack categories from two established source datasets, SALAD provides sufficient scale and diversity for training deep learning models while maintaining reproducibility through deterministic generation from public source data.

\item Difficulty levels (easy/medium/hard) and LLM-ready natural language descriptions open up new evaluation setups, such as difficulty-stratified benchmarking and prompt-based alert triage.

\item SALAD provides a reproducible public benchmark for RL/MARL-based alert triage research, where existing datasets are either proprietary, synthetic, or lack the multi-task evaluation structure needed for systematic comparison.
\end{itemize}


% ============================================================
\section{Data Description}
% ============================================================

\subsection{Dataset Overview}

SALAD contains 2,778,424 SOC alerts derived from two established network intrusion detection datasets: CIC-IDS2017~\cite{sharafaldin2018} (2,520,751 alerts) and UNSW-NB15~\cite{moustafa2015} (257,673 alerts). Table~\ref{tab:overview} summarizes the dataset.

\begin{table}[h]
\centering
\caption{SALAD Dataset Overview}
\label{tab:overview}
\begin{tabular}{lr}
\toprule
\textbf{Metric} & \textbf{Value} \\
\midrule
Total alerts & 2,778,424 \\
Columns (fields) & 29 \\
Attack categories & 15 \\
MITRE ATT\&CK tactics & 9 \\
MITRE ATT\&CK techniques & 12 \\
Kill chain phases & 8 \\
Malicious ratio & 21.2\% \\
Benign ratio & 78.8\% \\
Attack campaigns (correlated) & 586,227 \\
File size (full dataset) & 878 MB \\
\midrule
Training split & 1,944,887 (70\%) \\
Validation split & 416,756 (15\%) \\
Test split & 416,781 (15\%) \\
\bottomrule
\end{tabular}
\end{table}


\subsection{Schema Design}

Each alert record contains 29 fields organized into six groups (Table~\ref{tab:schema}).

\begin{table}[h]
\centering
\small
\caption{SALAD Schema (29 fields)}
\label{tab:schema}
\begin{tabular}{llp{6cm}}
\toprule
\textbf{Group} & \textbf{Field} & \textbf{Description} \\
\midrule
\multirow{3}{*}{Metadata} & \texttt{alert\_id} & Unique identifier \\
 & \texttt{timestamp} & ISO 8601 timestamp \\
 & \texttt{source\_dataset} & CICIDS2017 or UNSW\_NB15 \\
\midrule
\multirow{6}{*}{Network} & \texttt{src\_ip}, \texttt{dst\_ip} & Anonymized IP addresses \\
 & \texttt{src\_port}, \texttt{dst\_port} & Port numbers \\
 & \texttt{protocol} & TCP, UDP, etc. \\
 & \texttt{flow\_duration} & Duration in seconds \\
 & \texttt{total\_fwd/bwd\_packets} & Packet counts \\
 & \texttt{flow\_bytes\_per\_sec} & Throughput \\
\midrule
\multirow{5}{*}{Alert} & \texttt{alert\_type} & 16 alert categories \\
 & \texttt{severity} & none/low/medium/high/critical \\
 & \texttt{confidence} & Detection confidence [0,1] \\
 & \texttt{mitre\_tactic/technique} & MITRE ATT\&CK v14 mapping~\cite{strom2018mitre} \\
 & \texttt{kill\_chain\_phase} & Lockheed Martin kill chain + unknown~\cite{hutchins2011killchain} \\
\midrule
\multirow{4}{*}{Context} & \texttt{network\_segment} & dmz/server/workstation/mgmt \\
 & \texttt{alert\_count\_1h/24h} & Temporal alert frequency \\
 & \texttt{is\_repeated\_target} & Repeated target indicator \\
 & \texttt{correlation\_key} & Attack campaign identifier \\
\midrule
\multirow{4}{*}{Labels} & \texttt{is\_malicious} & Ground truth binary label \\
 & \texttt{attack\_category} & 15 attack categories \\
 & \texttt{triage\_decision} & escalate/investigate/suppress \\
 & \texttt{priority\_score} & Weighted priority [0,1] \\
\midrule
\multirow{2}{*}{LLM} & \texttt{difficulty\_level} & easy/medium/hard \\
 & \texttt{alert\_description} & Natural language description \\
\bottomrule
\end{tabular}
\end{table}


\subsection{Benchmark Tasks}

SALAD defines four evaluation tasks (Table~\ref{tab:tasks}).

\begin{table}[h]
\centering
\caption{SALAD Benchmark Tasks and Baseline Results}
\label{tab:tasks}
\begin{tabular}{llllc}
\toprule
\textbf{Task} & \textbf{Type} & \textbf{Target} & \textbf{Metric} & \textbf{RF Baseline} \\
\midrule
1. Alert Classification & Binary & \texttt{is\_malicious} & F1 & 0.9580 \\
2. Alert Triage & 3-class & \texttt{triage\_decision} & Weighted F1 & 0.9751 \\
3. Prioritization & Regression & \texttt{priority\_score} & Spearman & 0.7617 \\
4. Attack Category & Multi-class & \texttt{attack\_category} & Macro F1 & 0.5479 \\
\bottomrule
\end{tabular}
\end{table}


\subsection{Comparison with Existing Datasets}

Table~\ref{tab:comparison} compares SALAD with nine existing cybersecurity datasets~\cite{tavallaee2009kdd,sharafaldin2018,moustafa2015,msguide2024,landauer2024,husak2019,casilimas2024uwf,rodriguez2020mordor,sabu2023,cage2022cyborg} across 10 key dimensions.

\begin{table}[h]
\centering
\small
\caption{Comparison with Existing Datasets (\checkmark = fully supported, $\sim$ = partial, $\times$ = not available)}
\label{tab:comparison}
\begin{tabular}{lcccccccccc}
\toprule
\textbf{Dimension} & \rotatebox{90}{\textbf{KDD}} & \rotatebox{90}{\textbf{CIC-IDS}} & \rotatebox{90}{\textbf{UNSW}} & \rotatebox{90}{\textbf{MS GUIDE}} & \rotatebox{90}{\textbf{AIT-ADS}} & \rotatebox{90}{\textbf{SABU}} & \rotatebox{90}{\textbf{UWF-Zeek}} & \rotatebox{90}{\textbf{Mordor}} & \rotatebox{90}{\textbf{CybORG}} & \rotatebox{90}{\textbf{SALAD}} \\
\midrule
Public download        & \checkmark & \checkmark & \checkmark & \checkmark & \checkmark & \checkmark & \checkmark & \checkmark & \checkmark & \checkmark \\
Real traffic base      & \checkmark & \checkmark & $\sim$     & \checkmark & $\times$   & \checkmark & \checkmark & $\times$   & $\times$   & \checkmark \\
SOC alert format       & $\times$   & $\times$   & $\times$   & $\sim$     & $\sim$     & $\sim$     & $\times$   & $\times$   & $\times$   & \checkmark \\
Triage labels          & $\times$   & $\times$   & $\times$   & \checkmark & $\times$   & $\times$   & $\times$   & $\times$   & $\times$   & \checkmark \\
MITRE ATT\&CK          & $\times$   & $\times$   & $\times$   & $\times$   & $\times$   & $\times$   & \checkmark & \checkmark & $\sim$     & \checkmark \\
Unified schema         & \checkmark & \checkmark & \checkmark & \checkmark & $\times$   & $\times$   & \checkmark & $\times$   & N/A        & \checkmark \\
Multi-task benchmark   & $\times$   & $\times$   & $\times$   & $\times$   & $\times$   & $\times$   & $\times$   & $\times$   & \checkmark & \checkmark \\
RL/MARL ready          & $\times$   & $\times$   & $\times$   & $\times$   & $\times$   & $\times$   & $\times$   & $\times$   & \checkmark & \checkmark \\
LLM eval-ready         & $\times$   & $\times$   & $\times$   & $\times$   & $\times$   & $\times$   & $\times$   & $\times$   & $\times$   & \checkmark \\
Reproducible           & \checkmark & \checkmark & \checkmark & $\times$   & \checkmark & $\sim$     & \checkmark & \checkmark & \checkmark & \checkmark \\
\midrule
\textbf{Total \checkmark} & 4 & 4 & 4 & 4 & 2 & 2 & 5 & 3 & 4 & \textbf{10} \\
\bottomrule
\end{tabular}
\end{table}


% ============================================================
\section{Experimental Design, Materials, and Methods}
% ============================================================

\subsection{Source Datasets}

Two widely-cited, publicly available network intrusion detection datasets are used as source data:

\begin{enumerate}
\item \textbf{CIC-IDS2017}~\cite{sharafaldin2018}: Contains network traffic captured over five days (July 3--7, 2017) at the Canadian Institute for Cybersecurity. Includes 2,520,751 labeled flows across 14 attack types (grouped into 7 categories in the original paper).

\item \textbf{UNSW-NB15}~\cite{moustafa2015}: Contains 257,673 records combining normal and attack traffic across 9 attack categories, generated using the IXIA PerfectStorm tool at the University of New South Wales.
\end{enumerate}


\subsection{Conversion Pipeline}

The SALAD generation pipeline consists of five stages:

\textbf{Stage 1 --- Schema Mapping}: Each source record is mapped to the SALAD 29-field schema. Attack labels are mapped to standardized \texttt{alert\_type} categories, MITRE ATT\&CK tactics and techniques, and Cyber Kill Chain phases. IP addresses are anonymized using deterministic hashing for reproducibility. Severity levels are assigned based on attack type impact classification.

\textbf{Stage 2 --- Confidence Scoring}: Detection confidence scores are computed as: $c = \text{base}(m) + 0.3 \cdot s + 0.2 \cdot \min(1, f/10^6) + \epsilon$, where $m$ indicates malicious status (base = 0.3 for malicious, 0.15 for benign), $s$ is the severity score, $f$ is the flow bytes per second, and $\epsilon \sim U(-0.05, 0.05)$ adds realistic noise.

\textbf{Stage 3 --- Triage and Priority}: Triage decisions (escalate, investigate, suppress) are assigned using rule-based policies derived from SOC triage guidelines. Priority scores are computed as: $p = 0.4s + 0.3c + 0.2m + 0.1\min(1, n_{1h}/10)$, where $n_{1h}$ is the temporal alert count. Difficulty levels are assigned based on the combination of severity, confidence, and malicious status, producing easy (60.1\%), medium (33.0\%), and hard (6.9\%) distributions.

\textbf{Stage 4 --- Enrichment}: Temporal context features (\texttt{alert\_count\_1h}, \texttt{alert\_count\_24h}) are computed via per-source-IP rolling windows. Attack campaigns are identified using correlation keys based on source IP and kill chain phase combinations, yielding 586,227 unique campaigns. Natural language alert descriptions are generated from templates specific to each alert type.

\textbf{Stage 5 --- Splitting}: Stratified train (70\%) / validation (15\%) / test (15\%) splits are generated using joint stratification~\cite{sechidis2011stratification} by \texttt{attack\_category} and \texttt{source\_dataset} to ensure consistent class distributions across splits.

\subsubsection*{Label Validation}
The rule-based triage labels were validated through three checks: (1)~all escalate-labeled alerts correspond to high- or critical-severity attack types, matching standard SOC triage protocols; (2)~the class distribution (escalate~14.4\%, investigate~6.9\%, suppress~78.8\%) is consistent with reported real-world SOC alert distributions where 60--90\% of alerts are false positives or low-priority~\cite{alahmadi2022socsurvey,ponemon2023alert}; and (3)~the baseline Random Forest classifier achieves a weighted F1 of 0.9751 on triage prediction (Table~\ref{tab:results}), confirming that the labels encode learnable patterns from network-level features rather than random assignments.


\subsection{Baseline Evaluation}

We evaluate Random Forest~\cite{breiman2001rf} (100 trees) on all four tasks using only network-level features (ports, flow statistics, temporal counts) to avoid data leakage from label-derived features~\cite{kaufman2012leakage}. Table~\ref{tab:results} presents the detailed results.

\begin{table}[h]
\centering
\caption{Detailed Baseline Results --- Random Forest (100 trees), test set}
\label{tab:results}
\begin{tabular}{llcccc}
\toprule
\textbf{Task} & \textbf{Class/Metric} & \textbf{Precision} & \textbf{Recall} & \textbf{F1} & \textbf{Support} \\
\midrule
\multirow{3}{*}{1. Classification} & Benign & 0.98 & 0.99 & 0.99 & 328,211 \\
 & Malicious & 0.97 & 0.93 & 0.95 & 88,570 \\
 & \textit{Weighted avg} & \textit{0.96} & \textit{0.96} & \textit{0.96} & 416,781 \\
\midrule
\multirow{4}{*}{2. Triage} & Escalate & 0.95 & 0.94 & 0.94 & 59,993 \\
 & Investigate & 0.89 & 0.88 & 0.89 & 28,577 \\
 & Suppress & 0.99 & 0.99 & 0.99 & 328,211 \\
 & \textit{Weighted avg} & \textit{0.98} & \textit{0.98} & \textit{0.98} & 416,781 \\
\midrule
3. Prioritization & Spearman $\rho$ & \multicolumn{4}{c}{0.7617} \\
\midrule
4. Attack Category & Macro F1 & \multicolumn{4}{c}{0.5479 (15 classes)} \\
\bottomrule
\end{tabular}
\end{table}


\subsection{Reproducibility}

The complete pipeline is implemented in Python and publicly available at \url{https://github.com/nutthakorn7/SALAD-SOC}. Running \texttt{build\_salad\_fast.py} with the two source datasets (both publicly downloadable) deterministically reproduces the identical SALAD dataset. The build process completes in approximately 266 seconds on a standard workstation (Intel i7, 16~GB RAM). Baseline evaluation can be reproduced using \texttt{evaluate.py -{}-task all -{}-max-train 50000}.

\subsection{File Structure}

The released dataset contains the following files:

\begin{itemize}
\item \texttt{salad\_train.csv} --- Training split (1,944,887 alerts, 615~MB)
\item \texttt{salad\_val.csv} --- Validation split (416,756 alerts, 132~MB)
\item \texttt{salad\_test.csv} --- Test split (416,781 alerts, 132~MB)
\item \texttt{stats.json} --- Dataset statistics summary (JSON)
\item \texttt{datacard.md} --- Detailed dataset documentation
\item \texttt{evaluate.py} --- Evaluation and baseline scripts (Python)
\item \texttt{LICENSE} --- CC-BY 4.0 license
\end{itemize}


% ============================================================
\section{Ethics Statement}
% ============================================================

SALAD is derived entirely from two widely-used, publicly available datasets originally created for network security research (CIC-IDS2017 and UNSW-NB15). No new data collection involving human participants or animals was performed. All IP addresses in the source datasets are anonymized using deterministic hashing. No personally identifiable information (PII) is present in the released dataset. As a secondary analysis of existing public datasets, this work did not require institutional ethics review board approval. The dataset is released under a CC-BY 4.0 license.


% ============================================================
\section{Limitations}
% ============================================================

The dataset has several limitations that users should consider. First, SALAD alerts are synthetically derived from CIC-IDS2017 and UNSW-NB15 network traffic rather than captured from a live SOC environment, so the alert distribution and temporal patterns may not fully reflect production SOC workloads~\cite{ring2019ids_datasets}. Second, the triage decisions (escalate/investigate/suppress), severity scores, and priority labels are generated by deterministic rules based on attack type rather than labeled by human SOC analysts, which may introduce systematic biases absent from real analyst judgments. Third, the natural language alert descriptions are template-generated and lack the variability of vendor-specific alert formats. Finally, the dataset does not include contextual information such as asset criticality, network topology, or historical incident data that analysts typically use in triage decisions.


% ============================================================
\section{CRediT Author Statement}
% ============================================================

\textbf{Nutthakorn Chalaemwongwan}: Conceptualization, Methodology, Software, Data Curation, Validation, Writing -- Original Draft, Writing -- Review \& Editing.


% ============================================================
\section{Declaration of Competing Interest}
% ============================================================

The author declares that there are no known competing financial interests or personal relationships that could have appeared to influence the work reported in this paper.


% ============================================================
\section{Acknowledgements}
% ============================================================

This research received no external funding. During manuscript preparation the author used Claude (Anthropic) for English grammar checking and sentence-level editing. All research design, dataset construction, implementation, and interpretation of results were performed solely by the author.


\begin{thebibliography}{20}

% [1]-[10]: Background section (line 58)
\bibitem{alahmadi2022socsurvey}
B.A. Alahmadi, L. Axon, I. Martinovic, 99\% False positives: A qualitative study of SOC analysts' perspectives on security alarms, in: USENIX Security Symposium, 2022, pp. 2783--2800.

\bibitem{ponemon2023alert}
Vectra AI, 2023 State of Threat Detection: The Defender's Dilemma, Technical Report, Vectra AI, San Jose, CA, 2023.

\bibitem{husak2019}
M. Husak, M. Zadnik, V. Bartos, P. Sokol, Dataset of intrusion detection alerts from a sharing platform, Data in Brief 22 (2019) 929--933.

\bibitem{zuech2015ids_survey}
R. Zuech, T.M. Khoshgoftaar, R. Wald, Intrusion detection and big heterogeneous data: A survey, Journal of Big Data 2~(1) (2015) 3.

\bibitem{cage2022cyborg}
M. Kiely, M. Bowman, J. Skerman, C. Wang, CybORG: A Gym for the Development of Autonomous Cyber Agents, arXiv:2108.09118, 2022.

\bibitem{landauer2024}
M. Landauer, F. Skopik, M. Wurzenberger, Introducing a New Alert Data Set for Multi-Step Attack Analysis, in: Proceedings of the 17th Cyber Security Experimentation and Test Workshop (CSET), ACM, 2024.

\bibitem{msguide2024}
S. Mitra, R. Bhatt, et al., Microsoft GUIDE: A large-scale benchmark for security incident prediction, arXiv:2404.18710, 2024.

\bibitem{sharafaldin2018}
I. Sharafaldin, A.H. Lashkari, A.A. Ghorbani, Toward Generating a New Intrusion Detection Dataset and Intrusion Traffic Characterization, in: Proceedings of the 4th International Conference on Information Systems Security and Privacy (ICISSP), 2018, pp. 108--116.

\bibitem{moustafa2015}
N. Moustafa, J. Slay, UNSW-NB15: A comprehensive data set for network intrusion detection systems, in: Military Communications and Information Systems Conference (MilCIS), 2015, pp. 1--6.

\bibitem{strom2018mitre}
B.E. Strom, A. Applebaum, D.P. Miller, K.C. Nickels, A.G. Pennington, C.B. Thomas, MITRE ATT\&CK: Design and Philosophy, Technical Report, The MITRE Corporation, 2018.

% [11]: Value of the Data (line 68)
\bibitem{veeramachaneni2016aianalyst}
K. Veeramachaneni, I. Arnaldo, V. Korber, C. Bassias, K. Li, AI$^2$: Training a big data machine to defend, in: IEEE International Conference on Big Data Intelligence and Computing, 2016, pp. 49--56.

% [12]: Schema (line 141)
\bibitem{hutchins2011killchain}
E.M. Hutchins, M.J. Cloppert, R.M. Amin, Intelligence-driven computer network defense informed by analysis of adversary campaigns and intrusion kill chains, in: Leading Issues in Information Warfare \& Security Research, 2011, pp. 80--106.

% [13]-[16]: Comparison table (line 183)
\bibitem{tavallaee2009kdd}
M. Tavallaee, E. Bagheri, W. Lu, A.A. Ghorbani, A detailed analysis of the KDD CUP 99 data set, in: IEEE Symposium on Computational Intelligence for Security and Defense Applications (CISDA), 2009, pp. 1--6.

\bibitem{casilimas2024uwf}
K. Casilimas, J. Zhang, UWF-ZeekData24: A comprehensive network traffic dataset based on the CICIDS-2017 using Zeek, Data in Brief 56 (2024) 110813.

\bibitem{rodriguez2020mordor}
R. Rodriguez, J. Oltsik, Mordor Labs: Pre-recorded security events from simulated adversarial techniques, OTRF, 2020, \url{https://github.com/OTRF/Security-Datasets}.

\bibitem{sabu2023}
A. Sivanathan, H.H. Gharakheili, V. Sivaraman, SABU: A security alert benchmark for understanding network threat landscapes, IEEE Access 11 (2023) 85121--85135.

% [17]: Pipeline Stage 5 (line 238)
\bibitem{sechidis2011stratification}
K. Sechidis, G. Tsoumakas, I. Vlahavas, On the stratification of multi-label data, in: European Conference on Machine Learning and Knowledge Discovery in Databases (ECML PKDD), Springer, 2011, pp. 145--158.

% [18]-[19]: Baseline (line 243)
\bibitem{breiman2001rf}
L. Breiman, Random forests, Machine Learning 45~(1) (2001) 5--32.

\bibitem{kaufman2012leakage}
S. Kaufman, S. Rosset, C. Perlich, O. Stitelman, Leakage in data mining: Formulation, detection, and avoidance, ACM Transactions on Knowledge Discovery from Data 6~(4) (2012) 1--21.

% [20]: Limitations (line 300)
\bibitem{ring2019ids_datasets}
M. Ring, S. Wunderlich, D. Scheuring, D. Landes, A. Hotho, A survey of network-based intrusion detection data sets, Computers \& Security 86 (2019) 147--167.

\end{thebibliography}

\end{document}
